% ============================================================
% Section 169: 中心力场等效势的小振动近似
% ============================================================
\section{量子力学：中心力场等效势的小振动近似}
\label{sec:central-effective-potential-harmonic}

\begin{modelbox}[题目：中心力场等效势的小振动近似]
中心力场中等效势能 $V_l(r)=V(r)+\dfrac{l(l+1)\hbar^2}{2\mu r^2}$。对类氢离子（核电荷 $Ze$），$l\gg1$ 的情形：
\begin{enumerate}
    \item 求 $V_l(r)$ 的极小值及平衡距离 $r_l$；
    \item 将 $V_l(r)$ 在 $r_l$ 附近 Taylor 展开，保留 $(r-r_l)^2$ 项；
    \item 视电子径向运动为 $r_l$ 附近的小振动，求最低能级和激发能级，并与精确解比较。
\end{enumerate}
\end{modelbox}

\begin{tipbox}[考点快速分析]
\begin{itemize}
    \item \textbf{极小值条件}：$V_l'(r)=0$ 给出平衡距离
    \item \textbf{谐振近似}：Taylor 展开到二阶，化为标准谐振子方程
    \item \textbf{频率标度}：$\omega_l\propto r_l^{-3/2}\propto l^{-3}$
    \item \textbf{大 $l$ 极限}：小振动近似与精确氢原子能级一致
\end{itemize}
\end{tipbox}

\begin{derivationbox}[标准解题过程]
\textbf{1. 极小值与平衡距离}

类氢离子 $V(r)=-Ze^2/r$，等效势
\[
V_l(r)=\frac{l(l+1)\hbar^2}{2\mu r^2}-\frac{Ze^2}{r}.
\]

令 $V_l'(r)=0$：
\[
-\frac{l(l+1)\hbar^2}{\mu r^3}+\frac{Ze^2}{r^2}=0\quad\Rightarrow\quad r_l=\frac{l(l+1)\hbar^2}{\mu Ze^2}=\frac{l(l+1)a_0}{Z},
\]
其中 $a_0=\hbar^2/(\mu e^2)$。极小值
\[
V_l(r_l)=-\frac{Ze^2}{2r_l}=-\frac{Z^2e^2}{2l(l+1)a_0}.
\]

\textbf{2. Taylor 展开}

令 $r=r_l+\xi$，展开到二阶：
\[
V_l(r)\approx V_l(r_l)+\frac12V_l''(r_l)\xi^2.
\]

计算 $V_l''(r_l)$：
\[
V_l''(r)=\frac{3l(l+1)\hbar^2}{\mu r^4}-\frac{2Ze^2}{r^3},\quad V_l''(r_l)=\frac{Ze^2}{r_l^3}>0.
\]

于是 $V_l(r)\approx V_l(r_l)+\frac12\mu\omega_l^2\xi^2$，其中
\begin{equation}
\boxed{\omega_l=\sqrt{\frac{Ze^2}{\mu r_l^3}}=\frac{Z^2e^2}{\hbar a_0\,[l(l+1)]^{3/2}}.}
\end{equation}

\textbf{3. 小振动能级与精确解比较}

径向方程化为谐振子方程，能级
\begin{equation}
\boxed{E_{n_r,l}\approx V_l(r_l)+\Bigl(n_r+\frac12\Bigr)\hbar\omega_l,\quad n_r=0,1,2,\ldots}
\end{equation}

精确解 $E_n=-\dfrac{Z^2e^2}{2n^2a_0}$（$n=l+n_r+1$）。基态 ($n_r=0$) 精确能量 $E_{l+1}=-\dfrac{Z^2e^2}{2(l+1)^2a_0}$。小振动近似给出
\[
E_{0l}\approx-\frac{Z^2e^2}{2a_0}\Bigl[\frac{1}{l(l+1)}-\frac{1}{[l(l+1)]^{3/2}}\Bigr].
\]

$l\gg1$ 时展开：$\dfrac{1}{l(l+1)}\approx\dfrac{1}{l^2}$，$\dfrac{1}{[l(l+1)]^{3/2}}\approx\dfrac{1}{l^3}$，故
\[
E_{0l}\approx-\frac{Z^2e^2}{2a_0}\frac{1}{l^2}\Bigl(1-\frac{1}{l}\Bigr)\approx E_{l+1}+O\Bigl(\frac{1}{l^2}\Bigr).
\]

能级间隔：精确展开 $\Delta E_{\text{精确}}\approx\dfrac{Z^2e^2}{a_0(l+1)^3}$；小振动 $\hbar\omega_l\approx\dfrac{Z^2e^2}{a_0l^3}$。大 $l$ 下两者一致。

结论：小振动近似正确给出大 $l$ 时类氢离子能级的主要特征——准均匀分布，最低能级接近精确值。
\end{derivationbox}

\begin{formulabox}[答案汇总]
\begin{center}
\begin{tabular}{ll}
\toprule
\textbf{结论} & \textbf{结果} \\
\midrule
平衡距离 & $r_l=l(l+1)a_0/Z$ \\
极小势能 & $V_l(r_l)=-Z^2e^2/[2l(l+1)a_0]$ \\
振动频率 & $\omega_l=\sqrt{Ze^2/(\mu r_l^3)}$ \\
近似能级 & $E_{n_r,l}=V_l(r_l)+(n_r+\frac12)\hbar\omega_l$ \\
适用条件 & $l\gg1$，$n_r\ll l$ \\
\bottomrule
\end{tabular}
\end{center}
\end{formulabox}

\begin{insightbox}[物理图像]
\begin{itemize}
    \item 大 $l$ 时离心势垒很高，电子被``囚禁''在 $r_l$ 附近的环形区域内，径向运动幅度很小，谐振近似极其精确。
    \item 该图像是理解里德伯原子（高激发态氢原子）的基础：高 $l$ 的里德伯电子轨道接近经典圆形轨道，半径 $r_l\propto l^2$。
    \item 小振动近似失效于低 $l$ 或高 $n_r$（大径向振幅），此时需用精确类氢解或 WKB 法。
\end{itemize}
\end{insightbox}
